\documentclass[../../main.tex]{subfiles}
\begin{document}
\section{覆叠流形}\label{section:ch21-covering-manifolds}

命题 4.40 表明光滑流形的任意覆叠空间本身是光滑流形。了解被光滑流形覆盖的空间何时本身是光滑流形常很重要。为此需研究覆叠的自同构群在覆叠空间上的作用。第 7 章（命题 7.23）已看到，对任意光滑覆叠 $\pi\colon E\to M$，自同构群是离散 Lie 群，在覆叠空间 $E$ 上光滑自由作用。下面证明该作用也真。

先给出离散 Lie 群自由作用真的等价刻画。

\begin{lemma}\label{lemma:ch21-离散作用真判别}
设离散 Lie 群 $\Gamma$ 在流形 $E$ 上连续自由作用。该作用真当且仅当下列两条件都成立：(i) 每点 $p\in E$ 有邻域 $U$，使对每个 $g\in\Gamma$，只要 $g\neq e$ 就有 $(gU)\cap U=\varnothing$；(ii) 若 $p,q\in E$ 不在同一 $\Gamma$-轨道，则存在 $p$ 的邻域 $V$ 与 $q$ 的邻域 $W$，使对所有 $g\in\Gamma$ 有 $(gV)\cap W=\varnothing$。
\end{lemma}

\begin{proof}
先设作用自由且真，$\pi\colon E\to E/\Gamma$ 为商映射。由命题\cref{proposition:ch21-真作用Hausdorff}，$E/\Gamma$ Hausdorff。若 $p,q$ 不在同一轨道，可取 $\pi(p)$ 与 $\pi(q)$ 的不交邻域，其原像即满足 (ii)。为证 (i)，取 $p$ 的预紧邻域 $V$；由命题\cref{proposition:ch21-真作用刻画}，$\Gamma_V=\{g\colon(gV)\cap V\neq\varnothing\}$ 紧，而 $\Gamma$ 离散，故有限，写 $\Gamma_V=\{e,g_1,\dots,g_m\}$。必要时缩小 $V$ 使 $g_i^{-1}p\notin V$（即 $p\notin g_i V$）。则 $U=V\setminus(g_1 V\cup\cdots\cup g_m V)$ 满足 (i)。
反之，设 (i)(ii) 成立。设 $(g_i)\subset\Gamma$、$(p_i)\subset E$，$p_i\to p$、$g_i p_i\to q$。若 $p,q$ 不同轨道，由 (ii) 有邻域 $V,W$，但 $i$ 大时 $p_i\in V$、$g_i p_i\in W$，与 $(gV)\cap W=\varnothing$ 矛盾。故同轨道，存在 $g$ 使 $gp=q$，于是 $g^{-1}g_i p_i\to p$。取 (i) 中邻域 $U$，$i$ 大时 $p_i$ 与 $g^{-1}g_i p_i$ 都在 $U$；而 $(g^{-1}g_i U)\cap U\neq\varnothing$，由 (i) 得 $g^{-1}g_i=e$，即 $g_i=g$。故 $(g_i)$ 收敛，作用真。
\end{proof}

\begin{proposition}\label{proposition:ch21-自同构群作用真}
设 $M$ 是光滑流形，$\pi\colon E\to M$ 是光滑覆叠映射。赋予离散拓扑，自同构群 $\operatorname{Aut}(E)$ 在 $E$ 上光滑、自由、真作用。
\end{proposition}

\begin{proof}
命题 7.23 已证作用光滑自由。为证真，验证引理\cref{lemma:ch21-离散作用真判别} 的 (i)(ii)。任取 $p\in E$，取 $M$ 中 $\pi(p)$ 的均匀覆盖邻域 $W$；若 $U$ 为 $\pi^{-1}(W)$ 含 $p$ 的分支，则 (i) 易验。若 $p,q$ 不同轨道，同引理证明可取 $\pi(p)$、$\pi(q)$ 的不交邻域，其原像满足 (ii)。
\end{proof}

\begin{theorem}{[离散群的商为正规覆叠]}\label{theorem:ch21-离散群商覆叠}
设 $E$ 是连通光滑流形，$\Gamma$ 是在 $E$ 上光滑、自由、真作用的离散 Lie 群。则轨道空间 $E/\Gamma$ 是拓扑流形，并有唯一光滑结构使 $\pi\colon E\to E/\Gamma$ 是光滑正规覆叠映射。
\end{theorem}

\begin{proof}
由商流形定理，$E/\Gamma$ 有唯一光滑流形结构使 $\pi$ 是光滑淹没。光滑覆叠映射特别是光滑淹没，故使 $\pi$ 为光滑覆叠的任何其他光滑结构必等于此。由于 $\dim(E/\Gamma)=\dim E$，$\pi$ 是局部微分同胚。故只需证 $\pi$ 是覆叠映射。
取 $p\in E$。由引理\cref{lemma:ch21-离散作用真判别}，$p$ 有邻域 $U$ 使 $(gU)\cap U=\varnothing$（$g\neq e$）。必要时取连通。令 $V=\pi(U)$，由引理\cref{lemma:ch21-商映射开} 它在 $E/\Gamma$ 中开。因 $\pi^{-1}(V)$ 是互不相交连通开集 $gU$（$g\in\Gamma$）之并，且每个 $gU$ 上 $\pi$ 限制为同胚（用 $g$ 平移），故 $V$ 被均匀覆盖。元素 $g\in\Gamma$ 作为 $\pi$ 的自同构，按定义在纤维上可迁，故该覆叠正规。
\end{proof}

\begin{example}{[真离散群作用]}\label{example:ch21-真离散作用}
(a) 离散 Lie 群 $\mathbb{Z}^n$ 通过平移在 $\mathbb{R}^n$ 上光滑自由作用（例 7.22(e)）。若 $(x_i)$、$(m_i)$ 使 $x_i\to x$、$m_i+x_i\to y$，则 $m_i\to y-x$，由命题\cref{proposition:ch21-真作用刻画}(b) 作用真。轨道空间 $\mathbb{R}^n/\mathbb{Z}^n$ 同胚于 $n$-环面 $\mathbb{T}^n$；定理\cref{theorem:ch21-离散群商覆叠} 给出 $\mathbb{T}^n$ 上唯一使商映射为光滑覆叠的光滑结构。投影 $(x^1,\dots,x^n)\mapsto(e^{2\pi\mathrm{i}x^1},\dots,e^{2\pi\mathrm{i}x^n})$ 关于乘积光滑结构是局部微分同胚，且与商映射同一等同，故由定理 4.31，$\mathbb{R}^n/\mathbb{Z}^n\cong\mathbb{T}^n$。
(b) 二元群 $\{\pm1\}$ 通过数乘作用于 $\mathbb{S}^n$。该作用光滑自由，且群紧故真。这在 $\mathbb{S}^n/\{\pm1\}$ 上定义光滑结构。事实上该轨道空间微分同胚于 $\mathbb{RP}^n$：令 $q\colon\mathbb{S}^n\to\mathbb{RP}^n$ 为例 2.13(f) 的光滑覆叠，它与商映射 $\mathbb{S}^n\to\mathbb{S}^n/\{\pm1\}$ 同一等同，故由定理 4.31，$\mathbb{S}^n/\{\pm1\}\cong\mathbb{RP}^n$。
\end{example}
\end{document}
